\section{Results}
\label{sec:results}

We organise the results as six questions. Full per-set tables for every
architecture and metric are in Appendix~\ref{app:full-tables}; unless
stated otherwise, differences quoted below are paired Exact Match
differences with Holm-adjusted p-values, and ``significant'' means
$p<0.05$ after adjustment.

\subsection{RQ1: Is retrieval worth anything at all?}
\label{sec:results-worth}

Against the closed-book floor, naive retrieval is worth
\ClosedBookVsNaiveNQExactMatchDiff{} EM on NaturalQuestions,
\ClosedBookVsNaiveTriviaQAExactMatchDiff{} on TriviaQA,
\ClosedBookVsNaiveHotpotQAExactMatchDiff{} on HotpotQA, and
\ClosedBookVsNaiveMuSiQueExactMatchDiff{} on MuSiQue --- large,
significant gains everywhere except one set. On 2WikiMultihopQA the naive
row appears to \emph{lose} to no retrieval at all
(\ClosedBookVsNaiveTwoWikiExactMatchDiff{} EM), and this apparent
inversion is an abstention artefact, not a knowledge difference: holding
irrelevant passages, the reader declines on
\NaiveRAGTwoWikiAbstentionRate{} of questions (closed-book:
\ClosedBookTwoWikiAbstentionRate), and a decline scores zero where a
closed-book \emph{guess} is sometimes right. Restricted to questions both
systems answered, the ordering reverses
(\ClosedBookVsNaiveTwoWikiEMbothansweredDiff{} in naive's favour). Any
evaluation that reports EM without abstention would publish the wrong
conclusion on this set --- which is why every table in this article
carries both.

\subsection{RQ2: What does each single refinement buy?}
\label{sec:results-single}

\paragraph{Lexical and hybrid retrieval.} BM25 alone is well below the
dense baseline on the open-domain sets
(\NaiveVsBMTwentyFiveNQExactMatchDiff{} on NaturalQuestions) but
\emph{ahead} of it on 2WikiMultihopQA
(\NaiveVsBMTwentyFiveTwoWikiExactMatchDiff), whose questions name their
target articles verbatim --- exactly the rare-surface-form failure of
Section~\ref{sec:embeddings}. Fusing the two (Hybrid) is a small but
uniformly non-negative improvement
(\NaiveVsHybridNQExactMatchDiff{} to
\NaiveVsHybridTwoWikiExactMatchDiff{} across sets) at one CPU search of
extra cost --- the cheapest strictly-safe refinement measured here.

\paragraph{Cross-encoder reranking.} The strongest single gain and the
worst single failure, on different sets: \NaiveVsRerankHotpotQAExactMatchDiff{}
on HotpotQA and \NaiveVsRerankTwoWikiExactMatchDiff{} on 2WikiMultihopQA,
against \NaiveVsRerankNQExactMatchDiff{} on NaturalQuestions, where
answer-in-context falls from \NaiveRAGNQAnswerInContext{} to
\RerankNQAnswerInContext{} by the asking-for/asking-about mechanism of
Section~\ref{sec:cascade-fused}. Reranking is a bet on the question
distribution, not a safe default.

\paragraph{HyDE.} The one refinement that beats the baseline on
NaturalQuestions (\NaiveVsHyDENQExactMatchDiff{}, answer-in-context
\NaiveRAGNQAnswerInContext{} $\rightarrow$ \HyDENQAnswerInContext) and it
helps on the multi-hop sets as well
(\NaiveVsHyDETwoWikiExactMatchDiff{} on 2WikiMultihopQA), at the fixed
price of doubling the generator calls on every question.

\paragraph{Neighbour expansion.} Read against the twelve-passage control,
expansion is worth \NaiveTwelveVsNeighbourHotpotQAExactMatchDiff{} on
HotpotQA, \NaiveTwelveVsNeighbourTwoWikiExactMatchDiff{} on
2WikiMultihopQA and \NaiveTwelveVsNeighbourMuSiQueExactMatchDiff{} on
MuSiQue --- a null result, and a predicted one: the diagnosis
(RQ3) puts the ``wrong slice of the right article'' bucket this method
repairs at 5--14\% of failures. The popular advice to always fetch
neighbouring chunks is, on this corpus, worth approximately nothing once
context size is controlled.

\paragraph{Iterative and corrective methods.} Against the matched
ten-passage control, IRCoT gains
\NaiveTenVsIRCoTHotpotQAExactMatchDiff{} on HotpotQA,
\NaiveTenVsIRCoTTwoWikiExactMatchDiff{} on 2WikiMultihopQA and
\NaiveTenVsIRCoTMuSiQueExactMatchDiff{} on MuSiQue at 2.1--2.5 calls per
question; the one-shot variant with a worked example changes little
(\IRCoTVsIRCoTOneShotHotpotQAExactMatchDiff{} on HotpotQA,
\IRCoTVsIRCoTOneShotTwoWikiExactMatchDiff{} on 2WikiMultihopQA), because
the zero-shot loop already declares itself able to answer after the first
round --- the same five passages NaiveRAG reads --- on 70--80\% of
multi-hop questions. CRAG at the same budget gains
\NaiveTenVsCRAGTenMuSiQueExactMatchDiff{} on MuSiQue, its best set, at
three calls per question. Adaptive routing is indistinguishable from its
own single-hop branch (\NaiveVsAdaptiveTwoWikiExactMatchDiff{} vs.\ naive
on 2WikiMultihopQA): predicting difficulty from the question alone sends
92\% of a multi-hop set down the single-hop path.

\subsection{RQ3: Why does retrieval fail?}
\label{sec:results-diagnosis}

\begin{table}[t]
\centering\small
\caption{Diagnosis of the naive baseline's retrieval failures (300 sampled
failures per set; percentages of the diagnosed sample). The last row is
the Exact Match an oracle choosing per question between the naive run and
closed-book would add over the better single system.}
\label{tab:diagnosis}
\begin{tabular}{lrrr}
\toprule
Failure cause & HotpotQA & 2Wiki & MuSiQue \\
\midrule
Wrong slice of the right article & \DiagHotpotQAWrongSlicePct\% & \DiagTwoWikiWrongSlicePct\% & \DiagMuSiQueWrongSlicePct\% \\
Right article ranked 11--1000 & \DiagHotpotQARankedLowPct\% & \DiagTwoWikiRankedLowPct\% & \DiagMuSiQueRankedLowPct\% \\
Unreachable by this query & \DiagHotpotQAUnreachablePct\% & \DiagTwoWikiUnreachablePct\% & \DiagMuSiQueUnreachablePct\% \\
No gold article states the answer & \DiagHotpotQAAnswerAbsentPct\% & \DiagTwoWikiAnswerAbsentPct\% & \DiagMuSiQueAnswerAbsentPct\% \\
Gold article not in the corpus & \DiagHotpotQANotInCorpusPct\% & \DiagTwoWikiNotInCorpusPct\% & \DiagMuSiQueNotInCorpusPct\% \\
\midrule
Oracle (retrieve-or-closed-book) gain & \DiagHotpotQAOracleGain{} & \DiagTwoWikiOracleGain{} & \DiagMuSiQueOracleGain{} \\
\bottomrule
\end{tabular}
\end{table}

Table~\ref{tab:diagnosis} splits the baseline's failures into the five
causes of Section~\ref{sec:naive-pipeline}, each demanding a different
remedy. Three readings matter. First, the two \emph{unfixable} buckets ---
answer not stated by any gold article, article absent from the corpus ---
account for roughly a quarter to a third of failures, which is the
concrete form of the recall ceiling. Second, the \emph{ranking} bucket
(article present, reachable, ranked 11--1000) is the largest fixable one
on HotpotQA and 2WikiMultihopQA --- which is why reranking is the
strongest single refinement there --- while on MuSiQue the
\emph{unreachable} bucket dominates at \DiagMuSiQueUnreachablePct\%,
which is why query-side methods are what helps there. Third, the
\emph{wrong slice} bucket that neighbour expansion repairs is small
everywhere, predicting that method's null result. The diagnosis, made
before the cascade was designed, is the design brief the cascade answers:
fuse rankings (largest bucket), escalate to a query-side method
(second bucket), fall back to no retrieval (the questions where
retrieval cannot help but a guess might).

\subsection{RQ4: The cascade against every baseline}
\label{sec:results-main}

\begin{table}[t]
\centering\small
\caption{The proposed cascade against every baseline on every question
set. Each cell is the paired Exact Match difference (cascade minus
baseline); \textbf{bold} marks a significant cascade win and
\underline{underline} a significant loss under McNemar's test with Holm
adjustment over all \ResultPairs{} cells jointly. The cascade wins
\ResultWins{}, loses \ResultLosses{}, ties \ResultTies{}.}
\label{tab:main}
\ResultTable
\end{table}

Table~\ref{tab:main} is the headline result: \ResultWins{} significant
wins, \ResultLosses{} loss, \ResultTies{} ties over \ResultPairs{}
(baseline, question set) pairs, under the strictest correction. Against
the standard baseline the cascade gains
\NaiveVsCascadeTriviaQAExactMatchDiff{} on TriviaQA,
\NaiveVsCascadeHotpotQAExactMatchDiff{} on HotpotQA,
\NaiveVsCascadeTwoWikiExactMatchDiff{} on 2WikiMultihopQA and
\NaiveVsCascadeMuSiQueExactMatchDiff{} on MuSiQue; on NaturalQuestions the
difference (\NaiveVsCascadeNQExactMatchDiff) is not significant --- the
one set where the baseline was already near its measured ceiling for this
reader. The single loss is to reranking on HotpotQA
(\RerankVsCascadeHotpotQAExactMatchDiff), the set where the
cross-encoder's bet pays; against the same reranking on NaturalQuestions
the cascade gains \RerankVsCascadeNQExactMatchDiff. No other architecture
in the study has a row without a significant loss; several
(rerank, HyDE, CRAG-10) beat the cascade on their single best set. The
cascade's claim is not ``best everywhere'' but \emph{never worse than the
baseline, never catastrophic, at near-baseline cost} --- and
Table~\ref{tab:main} is that claim in test form.

Where does the gain come from? The stage distribution says: the fused
first stage answers \CascadeNQStageFused{} of \CascadeNQQuestions{}
NaturalQuestions questions, \CascadeTriviaQAStageFused{} of
\CascadeTriviaQAQuestions{} on TriviaQA, and
\CascadeMuSiQueStageFused{} of \CascadeMuSiQueQuestions{} on MuSiQue ---
the escalated remainder splits between HyDE
(e.g.\ \CascadeMuSiQueStageHyde{} on MuSiQue) and closed-book
(\CascadeMuSiQueStageClosedbook). The cascade's own abstention rate
collapses to \CascadeTriviaQAAbstentionRate--\CascadeMuSiQueAbstentionRate{}
(baseline: \NaiveRAGTriviaQAAbstentionRate--\NaiveRAGMuSiQueAbstentionRate),
because a question must be declined by \emph{three} strategies to end in a
refusal. Consistent with Section~\ref{sec:cascade-caveats}, EM restricted
to questions both systems answered is at or below zero against reranking
on four of five sets: the cascade wins by rescuing declined questions,
not by answering differently where its first stage already answered.

\subsection{RQ5: What does it cost?}
\label{sec:results-cost}

\begin{table}[t]
\centering\small
\caption{Cost per question, averaged over the sets each architecture ran
on, from the per-question logs. The last column is prompt tokens relative
to the naive baseline. The cross-encoder pass (rerank, fused, cascade) and
the CPU lexical search (BM25, hybrid, fused, cascade) are outside the
token count.}
\label{tab:cost}
\resizebox{\textwidth}{!}{\CostTable}
\end{table}

Table~\ref{tab:cost} prices every row. The cascade averages
\CostcascadeCalls{} generator calls and \CostcascadePromptRatio$\times$ the
baseline's prompt tokens --- against $2$--$3\times$ calls and
$2.2$--$4\times$ tokens for the corrective and iterative methods it
outperforms in Table~\ref{tab:main}. Figure~\ref{fig:pareto} plots both
dimensions: the cascade and its own first stage bound the quality--cost
frontier at each end, and every architecture that spends more than 1.6
calls per question sits strictly inside it. Escalation prices the
expensive strategies by the questions that need them: HyDE costs two
calls \emph{per escalated question}, and only 10--30\% of questions
escalate.

\subsection{RQ6: Which parts of the cascade matter?}
\label{sec:results-ablations}

\paragraph{What the fusion adds.} Stage one alone against plain reranking:
\RerankVsFusedNQExactMatchDiff{} on NaturalQuestions and
\RerankVsFusedHotpotQAExactMatchDiff{} on HotpotQA --- the fusion buys
back the collapse at the price of part of the gain, and its mean over the
five sets (0.322) is the best of any single-call architecture measured.
Swapping the first stage confirms the choice: a rerank-first cascade beats
the fused-first one on HotpotQA
(\CascadeRerankVsCascadeHotpotQAExactMatchDiff{} for fused-first) but
loses \CascadeRerankVsCascadeNQExactMatchDiff{} on NaturalQuestions; a
naive-first cascade is close behind everywhere
(\CascadeNaiveVsCascadeTwoWikiExactMatchDiff{} on 2WikiMultihopQA) and is
the cheapest variant; hybrid-first sits between
(\CascadeHybridVsCascadeNQExactMatchDiff{} on NaturalQuestions). No first
stage wins everywhere; fused has the best mean and the smallest
worst-case loss, which is its case.

\paragraph{What the escalation adds.} The cascade against its own first
stage run without escalation: \FusedVsCascadeNQExactMatchDiff{} on
NaturalQuestions, \FusedVsCascadeTriviaQAExactMatchDiff{} on TriviaQA,
\FusedVsCascadeHotpotQAExactMatchDiff{} on HotpotQA,
\FusedVsCascadeTwoWikiExactMatchDiff{} on 2WikiMultihopQA,
\FusedVsCascadeMuSiQueExactMatchDiff{} on MuSiQue --- one to three points
on every set, all significant, for 0.16--0.58 extra calls. The escalation
is worth roughly this much on top of \emph{any} first stage we tried;
which first stage is best depends on the set, but that the escalation
helps does not.

\paragraph{The confidence trigger does not pay.} The sweep over
\SweepQuestions{} held-out train questions puts the best threshold at
\SweepBestThreshold{} for a gain of $+0.0018$ EM over the abstention-only
cascade ($p=0.42$) at $+0.25$ calls per question. At that threshold on the
test sets, the trigger fires on 9--15\% of questions and the paired
difference against the abstention-only cascade is
\CascadeVsCascadeLogprobNQExactMatchDiff{} on NaturalQuestions,
\CascadeVsCascadeLogprobTriviaQAExactMatchDiff{} on TriviaQA,
\CascadeVsCascadeLogprobHotpotQAExactMatchDiff{} on HotpotQA,
\CascadeVsCascadeLogprobTwoWikiExactMatchDiff{} on 2WikiMultihopQA (the
one significant difference --- a loss) and
\CascadeVsCascadeLogprobMuSiQueExactMatchDiff{} on MuSiQue, at 0.2--0.36
extra calls. The greedy reader is nearly as confident when wrong as when
right: its token log-probabilities carry almost no information the
abstention did not already carry. We report this negative result because
the opposite is widely assumed; the abstention-only cascade remains the
proposed configuration. (Full sweep in Appendix~\ref{app:sweep}.)

\paragraph{The one-shot IRCoT objection.} Prompted with a worked example
as the original method prescribes, IRCoT changes by
\IRCoTVsIRCoTOneShotHotpotQAExactMatchDiff{} on HotpotQA and
\IRCoTVsIRCoTOneShotTwoWikiExactMatchDiff{} on 2WikiMultihopQA relative
to the zero-shot loop, and the cascade's comparison against it
(\IRCoTOneShotVsCascadeTwoWikiExactMatchDiff{} on 2WikiMultihopQA at
$2.4$ vs.\ $1.3$ calls) is materially the same as against the zero-shot
row. The multi-hop baseline was not under-prompted.

\subsection{RQ7: Does it survive a different reader?}
\label{sec:results-reader}

Everything the cascade claims rests on how a reader behaves when the
passages are irrelevant --- it must decline rather than guess. One model
is not evidence that readers in general do that, so the rows the claims
stand on (closed-book, naive, rerank, fused, cascade) were re-run with
Llama-3.1-8B-Instruct over the same index and questions, changing nothing
else. The pattern transfers: the cascade wins \ResultLlamaWins{} of
\ResultLlamaPairs{} pairs, loses \ResultLlamaLosses{}, ties
\ResultLlamaTies{} (Holm over the table), with the same signature ---
large gains over reranking on NaturalQuestions
(\RerankVsCascadeLlamaNQExactMatchDiff), a small loss to it on HotpotQA
(\RerankVsCascadeLlamaHotpotQAExactMatchDiff), significant gains over the
naive baseline on every set, and the escalation worth
\FusedVsCascadeLlamaNQExactMatchDiff{} to
\FusedVsCascadeLlamaTwoWikiExactMatchDiff{} over stage one alone. The
full second-reader table is in Appendix~\ref{app:full-tables}.
